% =====================================================================
%  sections/05_system_design.tex  --  V. SYSTEM DESIGN
%  Each subsection holds its diagram and nothing else.
% =====================================================================
\clearpage
\section{SYSTEM DESIGN}
\label{sec:system-design}

% ---------------------------------------------------------------------
\subsection{System Architecture Diagram}
\label{sub:system-architecture}

The diagram below presents a high-level view of how \projname{}'s major
layers fit together, from the client-facing web application down to the
data and infrastructure that support it. It traces how a user request
travels from the Next.js frontend through the FastAPI backend services
to the PostgreSQL and analytics data layer, illustrating the boundaries
of responsibility between presentation, business logic, and storage.

\projfigure{0.95}{fig-system-architecture}{\protect\projnameplain{}'s overall system architecture, showing the flow between users, the Next.js frontend, the FastAPI backend services, the PostgreSQL and analytics data layer, and the underlying infrastructure components.}{fig:system-architecture}

% ---------------------------------------------------------------------
\clearpage
\subsection{AI Architecture Diagram}
\label{sub:ai-architecture}

This diagram zooms into the AI layer introduced above, showing how a
natural-language question is turned into a validated SQL query and a
readable answer. It lays out the eight-step LangGraph workflow, from
the input guardrail and language detection through schema retrieval,
SQL generation, validation, and execution, ending with insight
generation for the user.

\projfigure{0.95}{fig-ai-architecture}{\protect\projnameplain{}'s AI agent pipeline, detailing the eight-step LangGraph workflow from input guardrail and language detection through schema retrieval, SQL generation, validation, execution, and insight generation.}{fig:ai-architecture}

% ---------------------------------------------------------------------
\clearpage
\subsection{Use Case Diagram}
\label{sub:use-case-diagram}

The use case diagram summarizes what each type of user can do within
\projname{} from a functional standpoint, independent of the
underlying architecture. It separates the actions available to the two
system roles, Admin and User, covering report management, chatbot
interaction, forecasting, and account administration.

\projfigure{0.95}{fig-use-case-diagram}{Use case diagram showing the distinct actions available to Admin and User roles, including report management, chatbot interaction, forecasting, and account administration.}{fig:use-case}

% ---------------------------------------------------------------------
\clearpage
\subsection{Layout Design}
\label{sub:layout-design}

The following figures walk through \projname{}'s main screens in the
order a typical user would encounter them: signing in, viewing the
dashboard, conversing with the AI chatbot, working with generated
reports, creating and reviewing forecasts, and, for administrators,
managing users and account settings. Together they illustrate the
platform's overall visual language and how each core feature is
presented to the end user.

\projfigure{0.90}{fig-layout-login}{\protect\projnameplain{}'s login page, showing the email and password authentication form alongside the platform's introductory branding panel.}{fig:layout-login}
\clearpage

\projfigure{0.90}{fig-layout-dashboard}{\protect\projnameplain{}'s main dashboard, displaying key financial KPIs, sales versus purchases trends, revenue breakdown by transaction type, and profit composition for the selected period.}{fig:layout-dashboard}
\projfigure{0.90}{fig-layout-chatbot-interface}{\protect\projnameplain{}'s AI chatbot landing view, showing suggested starter questions and the natural-language query input used to begin a conversation.}{fig:layout-chatbot-interface}
\clearpage

\projfigure{0.90}{fig-layout-chatbot-response}{Example chatbot response showing a natural-language insight, a generated bar chart visualization, and suggested follow-up questions for continued exploration.}{fig:layout-chatbot-response}
\projfigure{0.90}{fig-layout-all-reports}{Report management page listing all previously generated reports, with an option to create a new customizable report.}{fig:layout-all-reports}
\clearpage

\projfigure{0.90}{fig-layout-report-page}{Example generated financial report, showing KPI summaries, revenue versus spending trends, payment status, and settlement versus receivables visualizations.}{fig:layout-report-page}
\projfigure{0.90}{fig-layout-create-forecast}{Forecast creation page shown before any forecast has been generated, prompting the user to initiate a sales or purchase forecast.}{fig:layout-create-forecast}
\clearpage

\projfigure{0.90}{fig-layout-forecast}{Generated sales forecast, showing predicted values with confidence intervals, model performance metrics, and the key drivers influencing the forecast.}{fig:layout-forecast}
\projfigure{0.90}{fig-layout-user-management}{Admin-only user management page, listing all system users with their status, creation date, and available account actions.}{fig:layout-user-management}
\clearpage

\projfigure{0.90}{fig-layout-profile}{User profile page, showing editable personal information and account security settings such as password and recovery email.}{fig:layout-profile}